%-----------------------------------------------------------------------
% Title, authors, abstract, keywords
%-----------------------------------------------------------------------
\title[Wave-Corrected Caustics over Photon Splatting]%
      {Wave-Corrected Caustics over Photon Splatting:\\[2pt]
       Fold Kinematics, a Band-Local Airy Correction,\\ and a Comparative Evaluation Against Production Solvers}

% --- For anonymous review, replace this block with the submission ID.
\author[Debyte]
{\parbox{\textwidth}{\centering Debyte\thanks{Corresponding author.}\orcid{0000-0000-0000-0000}
        }
        \\
{\parbox{\textwidth}{\centering $^1$Independent research
       }
}}

%--- abstract -----------------------------------------------------------
\begin{abstract}
Photon-splatting and path-traced caustics are the standard way to render the
light focused by a curved transparent body onto a diffuse receiver. Both
families evaluate a geometric push-forward whose Jacobian determinant vanishes
on the fold locus, so the estimator diverges there, and no amount of sampling
repairs it. We state that failure as a property of the model rather than of the
discretisation, then describe a hybrid estimator in which a geometric quadrature
carries the whole image and a wave-optical correction is applied only inside a
narrowly detected fold band: the band's own geometric energy is redistributed
according to the Airy fold profile $|\Ai(-s/w)|^2$, with a per-channel width
$w_c$ following the $\lambda^{2/3}$ scaling implied by the fold's canonical
form. We first derive the kinematics from scratch --- the ray map, the fold
condition, the normal form of the fold and the wavelength scaling that follows
from it --- and then use that derivation to explain where the standard computer
graphics constructions for caustics succeed and where they necessarily fail.

We evaluate against PBRT-v4 SPPM, Blender Cycles and Mitsuba~3 on an identical
scene with measured dispersive indices (water: Daimon \& Masumura 2007; glass:
Rubin 1985). Against an independent brute-force geometric ground truth of
$810{,}000$ explicit rays, our estimator scores $0.917$ correlation at 40\,mm
tiles against $0.853$ (SPPM) and $0.828$ (Cycles); at 5\,mm tiles the ordering
is $0.721$ / $0.722$ / $0.405$. The converged two-dimensional shape of our field
agrees with SPPM at $r=+0.932$, and its transport throughput measures $96.41\%$
at three wavelengths with Fresnel reflection as the only loss channel.

We are explicit about what the method does \emph{not} establish. The correction
is defined only for a fold; it says nothing about higher-order caustics, and we
show by direct measurement that a naive normalisation of the band is
\emph{impossible} --- $|\Ai(-u)|^2$ is not integrable on $\mathbb{R}$ --- which
is why the correction is written as a discrete band ratio rather than a
continuous constant. Two of our own earlier claims are withdrawn in this paper,
each killed by an estimator artefact rather than by the physics, and both are
documented because they are the failure mode we expect a reader to hit.
\end{abstract}

%--- keywords -----------------------------------------------------------
\keywords{caustics, wave optics, photon splatting, photon mapping, SPPM,
diffraction, Airy function, grid-free rendering, comparative evaluation}
